\section{Wykaz oznaczeń}

\subsection*{Symbole matematyczne}

\begin{tabular}{@{}p{2.5cm}p{11.75cm}@{}}
  $\mathbf{B}$                                   & wektor bistyczny (ang. \textit{bitangent}), prostopadły do wektora stycznego i normalnego, wykorzystywany w macierzy TBN                                       \\[0.5em]
  $CV$                                           & współczynnik zmienności (ang. \textit{Coefficient of Variation}), miara względnego rozproszenia wyników, wyrażona wzorem $CV = \frac{\sigma}{\mu} \cdot 100\%$ \\[0.5em]
  $\mathbf{E}_1$, $\mathbf{E}_2$                 & wektory krawędzi trójkąta, które są definiowane jako $\mathbf{E}_1 = \mathbf{P}_1 - \mathbf{P}_0$ oraz $\mathbf{E}_2 = \mathbf{P}_2 - \mathbf{P}_0$            \\[0.5em]
  $k_d$                                          & współczynnik odbicia rozproszonego (ang. \textit{diffuse}), określający intensywność światła rozproszonego                                                     \\[0.5em]
  $k_s$                                          & współczynnik odbicia lustrzanego (ang. \textit{specular}), określający intensywność refleksów świetlnych                                                       \\[0.5em]
  $\mu$                                          & średnia arytmetyczna próby pomiarowej                                                                                                                          \\[0.5em]
  $\mathbf{N}$                                   & wektor normalny do powierzchni, prostopadły do płaszczyzny trójkąta                                                                                            \\[0.5em]
  $\mathbf{P}_0$, $\mathbf{P}_1$, $\mathbf{P}_2$ & wierzchołki trójkąta w przestrzeni trójwymiarowej                                                                                                              \\[0.5em]
  $\sigma$                                       & odchylenie standardowe, miara rozproszenia wartości wokół średniej                                                                                             \\[0.5em]
  $\mathbf{T}$                                   & wektor styczny (ang. \textit{tangent}), leżący w płaszczyźnie powierzchni i wskazujący kierunek wzrostu współrzędnej $u$ tekstury                              \\[0.5em]
  $(u, v)$                                       & współrzędne tekstury (ang. \textit{UV coordinates}), określające położenie punktu na dwuwymiarowej mapie tekstury                                              \\[0.5em]
  $\Delta u_1$, $\Delta u_2$                     & różnice współrzędnych tekstury $u$ między wierzchołkami trójkąta                                                                                               \\[0.5em]
  $\Delta v_1$, $\Delta v_2$                     & różnice współrzędnych tekstury $v$ między wierzchołkami trójkąta                                                                                               \\[0.5em]
\end{tabular}

\subsection*{Terminologia techniczna}

\begin{tabular}{@{}p{3.5cm}p{10.75cm}@{}}
  Renderowanie & proces przetwarzania danych przez specjalny potok (nazywany potokiem renderowania) przy użyciu procesora graficznego w celu ich przedstawienia na ekranie. \\[0.5em]

  Shader       & program wykonywany na procesorze graficznym.                                                                                                               \\[0.5em]
\end{tabular}

\newpage